% ============================================================================
%  代码图（Figure，library-paper 风格，对标 LLM Reasoners 的 code-to-class 图）
%  依赖：在 iclr2026_conference.tex 的 preamble 里加入下面 PREAMBLE 块（listings + xcolor 已加载）。
%  用法：% ============================================================================
%  代码图（Figure，library-paper 风格，对标 LLM Reasoners 的 code-to-class 图）
%  依赖：在 iclr2026_conference.tex 的 preamble 里加入下面 PREAMBLE 块（listings + xcolor 已加载）。
%  用法：% ============================================================================
%  代码图（Figure，library-paper 风格，对标 LLM Reasoners 的 code-to-class 图）
%  依赖：在 iclr2026_conference.tex 的 preamble 里加入下面 PREAMBLE 块（listings + xcolor 已加载）。
%  用法：% ============================================================================
%  代码图（Figure，library-paper 风格，对标 LLM Reasoners 的 code-to-class 图）
%  依赖：在 iclr2026_conference.tex 的 preamble 里加入下面 PREAMBLE 块（listings + xcolor 已加载）。
%  用法：\input{fig_code} 放在 intro 或 §3.2 Library Design 附近。
% ============================================================================

\begin{figure}[t]
\centering
\begin{lstlisting}[style=llmrouter]
from llmrouter.models import MetaRouter, BaseTrainer

class MyRouter(MetaRouter):                    # (E_q, E_m, g, d): state -> action
    def route_single(self, query):
        s = self.encode_state(query)           # context encoder  E_q
        scores = self.score(s, self.models)    # model encoder E_m + scoring g
        query["model_name"] = self.decide(scores)   # decision rule d
        return query

class MyRouterTrainer(BaseTrainer):            # learning signal  L
    def loss_func(self, outputs, batch):
        return my_objective(outputs, batch)    # pointwise / pairwise / RL reward

# Train and run through the same interface as every built-in router.
router  = MyRouter(yaml_path="my_router.yaml")
trainer = MyRouterTrainer(router)
trainer.train()
answer  = router.route_single({"query": "..."})
\end{lstlisting}
% \vspace{-4pt}
\caption{\textbf{The five components of the routing formulation map onto two classes in \method}. A router subclasses \texttt{MetaRouter} and implements \texttt{route\_single} (or \texttt{route\_batch}), where the context encoder $E_q$, model encoder $E_m$, scoring function $g$, and decision rule $d$ turn a state into a selected model; the learning signal $\mathcal{L}$ lives in a \texttt{BaseTrainer} subclass. }
\label{fig:code}
\vspace{-10pt}
\end{figure}
 放在 intro 或 §3.2 Library Design 附近。
% ============================================================================

\begin{figure}[t]
\centering
\begin{lstlisting}[style=llmrouter]
from llmrouter.models import MetaRouter, BaseTrainer

class MyRouter(MetaRouter):                    # (E_q, E_m, g, d): state -> action
    def route_single(self, query):
        s = self.encode_state(query)           # context encoder  E_q
        scores = self.score(s, self.models)    # model encoder E_m + scoring g
        query["model_name"] = self.decide(scores)   # decision rule d
        return query

class MyRouterTrainer(BaseTrainer):            # learning signal  L
    def loss_func(self, outputs, batch):
        return my_objective(outputs, batch)    # pointwise / pairwise / RL reward

# Train and run through the same interface as every built-in router.
router  = MyRouter(yaml_path="my_router.yaml")
trainer = MyRouterTrainer(router)
trainer.train()
answer  = router.route_single({"query": "..."})
\end{lstlisting}
% \vspace{-4pt}
\caption{\textbf{The five components of the routing formulation map onto two classes in \method}. A router subclasses \texttt{MetaRouter} and implements \texttt{route\_single} (or \texttt{route\_batch}), where the context encoder $E_q$, model encoder $E_m$, scoring function $g$, and decision rule $d$ turn a state into a selected model; the learning signal $\mathcal{L}$ lives in a \texttt{BaseTrainer} subclass. }
\label{fig:code}
\vspace{-10pt}
\end{figure}
 放在 intro 或 §3.2 Library Design 附近。
% ============================================================================

\begin{figure}[t]
\centering
\begin{lstlisting}[style=llmrouter]
from llmrouter.models import MetaRouter, BaseTrainer

class MyRouter(MetaRouter):                    # (E_q, E_m, g, d): state -> action
    def route_single(self, query):
        s = self.encode_state(query)           # context encoder  E_q
        scores = self.score(s, self.models)    # model encoder E_m + scoring g
        query["model_name"] = self.decide(scores)   # decision rule d
        return query

class MyRouterTrainer(BaseTrainer):            # learning signal  L
    def loss_func(self, outputs, batch):
        return my_objective(outputs, batch)    # pointwise / pairwise / RL reward

# Train and run through the same interface as every built-in router.
router  = MyRouter(yaml_path="my_router.yaml")
trainer = MyRouterTrainer(router)
trainer.train()
answer  = router.route_single({"query": "..."})
\end{lstlisting}
% \vspace{-4pt}
\caption{\textbf{The five components of the routing formulation map onto two classes in \method}. A router subclasses \texttt{MetaRouter} and implements \texttt{route\_single} (or \texttt{route\_batch}), where the context encoder $E_q$, model encoder $E_m$, scoring function $g$, and decision rule $d$ turn a state into a selected model; the learning signal $\mathcal{L}$ lives in a \texttt{BaseTrainer} subclass. }
\label{fig:code}
\vspace{-10pt}
\end{figure}
 放在 intro 或 §3.2 Library Design 附近。
% ============================================================================

\begin{figure}[t]
\centering
\begin{lstlisting}[style=llmrouter]
from llmrouter.models import MetaRouter, BaseTrainer

class MyRouter(MetaRouter):                    # (E_q, E_m, g, d): state -> action
    def route_single(self, query):
        s = self.encode_state(query)           # context encoder  E_q
        scores = self.score(s, self.models)    # model encoder E_m + scoring g
        query["model_name"] = self.decide(scores)   # decision rule d
        return query

class MyRouterTrainer(BaseTrainer):            # learning signal  L
    def loss_func(self, outputs, batch):
        return my_objective(outputs, batch)    # pointwise / pairwise / RL reward

# Train and run through the same interface as every built-in router.
router  = MyRouter(yaml_path="my_router.yaml")
trainer = MyRouterTrainer(router)
trainer.train()
answer  = router.route_single({"query": "..."})
\end{lstlisting}
% \vspace{-4pt}
\caption{\textbf{The five components of the routing formulation map onto two classes in \method}. A router subclasses \texttt{MetaRouter} and implements \texttt{route\_single} (or \texttt{route\_batch}), where the context encoder $E_q$, model encoder $E_m$, scoring function $g$, and decision rule $d$ turn a state into a selected model; the learning signal $\mathcal{L}$ lives in a \texttt{BaseTrainer} subclass. }
\label{fig:code}
\vspace{-10pt}
\end{figure}
